\chapter{AIGC 与迷因生产链：变异引擎的成本革命}\label{ch:08}

\begin{nlkeypoints}
  \begin{enumerate}[label=(\arabic*),leftmargin=2.4em]
    \item 奶蛙二创的生产链可拆为五环节（素材获取、形象生成、
          视频生成、音轨处理、分发），其中前四环节在 2024—2026 年间
          被 AIGC 工具相继改造；\keyw{变异引擎}是形象的 AI 重绘与
          视频生成模型（据二手评述，抖音 \tsm{Seedance} 于 2026 年
          上线后扩散显著加速，\conf{C} 级）。
    \item AIGC 给了第六章“一软两硬”结构一个经济学解释：
          \keyw{内容变异的边际成本趋近于零，形式构件作为模板
          被工具固化}，均衡态因此是“内容极活泼、形式极保守”
          （本书分析，\conf{M} 级）。
    \item 本章同时给出全书第一个\keyw{可复现方法}的完整样例：
          空耳词等价项的三条判定标准与“拼音声韵母编辑距离”
          排序指标（\conf{M} 级，仅为排序用），
          并以六个衍生语的空耳变体作演示计算。
  \end{enumerate}
\end{nlkeypoints}

前三章（符号、机制）都留下过同一类承诺：
“为什么奶蛙在 2026 年才爆发”“为什么形式稳定而内容狂变”，
答案都指向生产端。本章的任务是把生产端摊开：
谁在生产、用什么工具、每个环节的成本发生了什么变化、
以及——作为全书方法论承诺的兑现——其中一类语言材料
（空耳词）的判定如何做成\keyw{可复现的操作}。
本章的宏观背景口径（短视频用户规模等）以中国互联网络
信息中心的历次统计报告为准\cite{baidu-index}，
本书不引用具体数值，读者引用前应核对当期报告。

\section{生产链拆解：五个环节}\label{sec:08:chain}

把一条典型奶蛙二创视频倒放一遍，可以拆出五个生产环节。
表 \ref{tab:08:chain} 给出各环节的传统做法、AIGC 改造与
成本变化方向（拆解为本书分析，\conf{M} 级；
工具描述的来源均为 \conf{C} 级二手评述
\cite{shoutanji2026naiwa} 与社群考据
\cite{lostmedia2026naiwa}）。

\begin{table}[htbp]\centering
  \caption{奶蛙二创生产链的五个环节（本书拆解，\conf{M} 级）}
  \label{tab:08:chain}
  \small
  \begin{tabularx}{\textwidth}{@{}llXX@{}}
    \toprule
    环节 & 传统做法 & AIGC 改造 & 成本变化 \\ \midrule
    1 素材获取 & 手工截图、抠图 &
      智能体平台生成角色形象（如星野 AI 的头像生成） &
      下降：从“找图”变为“生成图” \\
    2 形象生成 & 手工改图（换皮、P 图） &
      图像模型按提示词批量出变体 &
      大幅下降：变异供给近乎无限 \\
    3 视频生成 & 逐帧动画、剪辑拼接 &
      视频生成模型直接出动态片段 &
      大幅下降：静态形象获得“演技” \\
    4 音轨处理 & 手动变调、对轨 &
      变调已成模板化操作 &
      下降后企稳：音轨本身标准化 \\
    5 分发 & 论坛、群聊 &
      短视频平台算法推荐 &
      结构变化：分发不再是瓶颈 \\
    \bottomrule
  \end{tabularx}
\end{table}

\begin{figure}[htbp]\centering
\begin{tikzpicture}[font=\small,
  node distance=0.45cm,
  st/.style={rounded corners,draw=NLDeep,thick,fill=NLYellow!18,
    align=center,minimum height=1.0cm,minimum width=2.35cm},
  tool/.style={font=\footnotesize,color=NLAmber,align=center},
  arr/.style={-{Stealth[length=2.6mm]},thick,NLDeep}]
  \node[st] (s1) {素材获取};
  \node[st,right=of s1] (s2) {形象生成};
  \node[st,right=of s2] (s3) {视频生成};
  \node[st,right=of s3] (s4) {音轨处理};
  \node[st,right=of s4] (s5) {分发};
  \node[tool,below=0.30cm of s1] {智能体平台\\头像生成};
  \node[tool,below=0.30cm of s2] {图像模型\\批量变体};
  \node[tool,below=0.30cm of s3] {视频生成模型\\（Seedance 等）};
  \node[tool,below=0.30cm of s4] {变调模板};
  \node[tool,below=0.30cm of s5] {算法推荐};
  \draw[arr] (s1) -- (s2);
  \draw[arr] (s2) -- (s3);
  \draw[arr] (s3) -- (s4);
  \draw[arr] (s4) -- (s5);
  \draw[arr,dashed,NLAmber] (s5.south) to[out=250,in=290]
    node[below,font=\footnotesize]{反馈：完播 / 转发数据反哺变异方向} (s2.south);
\end{tikzpicture}
  \caption{奶蛙二创生产链与 AIGC 工具的接入点（\conf{M} 级拆解；
          虚线为平台数据对上游变异方向的反馈）。}
  \label{fig:08:chain}
\end{figure}

\begin{nlhistory}[史料与考据：变异引擎的两个时点]
  \textbf{2024-10-02}：星野 AI 用户“纲手【百豪】”创建“奶龙”智能体，
  平台对用户输入图片的 AI 重绘产生了最初奶蛙形象——
  生产链环节 1 的意外产物成为整个迷因的起点
  （考据，\conf{C}\cite{lostmedia2026naiwa}；见第四章）。
  \textbf{2026 年上半年}：抖音视频生成模型 \tsm{Seedance} 推出，
  奶蛙视频的生产门槛显著降低，被二手评述列为
  奶蛙爆发期的供给侧变量
  （\conf{C}\cite{shoutanji2026naiwa}）。
  两个时点之间的 2024—2025 年，恰是第四章所述的
  “低可见度潜伏期”——\keyw{形象已变异、动态未解锁}，
  与第五章命题 \ref{prop:05:second}（二次爆发需要
  足够强的重新激活项）在时间线上相容。
\end{nlhistory}

必须声明这条时间线的证据地位：两个时点均为 \conf{C} 级，
且“Seedance 促成爆发”只是评述者的归因，
不是受控结论。它在本书中的用法是
\keyw{为 H7.5 的检验提供须控制的供给侧变量清单}
（第七章 \ref{sec:07:reverse} 的证伪条件之一），
而非作为已证事实参与论证。

\section{成本结构：形式稳定性的经济学版本}\label{sec:08:cost}

第六章命题 \ref{prop:06:formstable} 给出了形式稳定的
三重来源（学习成本、识别成本、组合开放）。
生产端视角可以把它压缩成一句更锋利的话：

\begin{proposition}[成本不对称命题（教学命题）]\label{prop:08:asym}
  当图像与视频生成把\keyw{内容变异}的边际成本压到趋近于零，
  而平台分发以\keyw{识别速度}为筛选条件时，
  稳定的形式构件（固定音轨、结尾套路、衍生语）
  成为唯一不可替代的资产：
  内容层发生“无限供给下的激烈淘汰”，
  形式层发生“模板固化”——均衡态即
  “内容极活泼、形式极保守”。
\end{proposition}

这条命题（\conf{M} 级）给出两个可检验推论。
推论一：AIGC 工具普及后，单条奶蛙视频的
\keyw{内容新颖度}上升而\keyw{形式新颖度}不升——
可通过内容编码（场景、题材的多样性指数）
与形式编码（音轨、结构的多样性指数）的对比检验。
推论二：形式构件的变异失败率应高于内容变异
——“改结尾套路”的二创传播表现应系统性差于
“同结构换内容”的二创，因为前者破坏了识别信号。
两个推论的编码方案第十一章一并给出。

反过来，这条命题也解释了第六章“窄通道规律”
（命题 \ref{prop:06:channel}）的一个竞争性解释：
轮廓趋同可能不是社群选择的结果，而是
\keyw{工具模板}的产物——图像模型对同一提示词的
输出分布本来就集中。判别方法在第六章已给出
（工具固定/社群固定的对照样本），此处只需指出：
生产端证据越强，窄通道规律归因于社群选择的
置信度就要相应调低。这是竞争性解释之间
\keyw{此消彼长}的典型形态。

\begin{nlcritique}[批判性反思：低质内容的规模问题]
  生产成本趋零的另一面是\keyw{供给的泛滥}。
  《人民日报》报道明确将“用 AI 工具批量生成低质内容”
  列为奶龙 IP 遭遇的侵权形态之一
  \conf{A}\cite{jiang2026nailong}——这是全书与
  AIGC 直接相关的最强的一条 \conf{A} 级材料，
  它说明“批量生成”不是研究者的抽象担忧，
  而是权利人实际维权清单上的项目。
  从传播政治经济学看，这里的结构性问题是：
  \keyw{生产成本由个体承担（趋零），拥挤成本由
  公共信息流承担（趋增）}——每个二创作者的最优策略
  （多产、快产）加总起来是全体用户的信息流污染。
  第六章的审丑批判、第十章的校园刷屏争议，
  都可以视为这一外部性的不同侧面。
  本书不评判任何具体二创行为的合法性（第十二章），
  但要求读者带着这个外部性结构去理解
  官方、平台与社群三方各自的反应。
\end{nlcritique}

\section{可复现方法：空耳词的判定与度量}\label{sec:08:soramimi}

第六章表 \ref{tab:06:langs} 列出了六个衍生语及其
“等效空耳”。本节兑现方法论承诺：把“等效”
变成可执行的判定与可复现的度量。

\subsection{三条判定标准}

沿用事实手册的口径，一个候选空耳变体 $w'$ 与
衍生语 $w$ 建立\keyw{等价关系}，须同时满足：

\begin{enumerate}[label=(\arabic*),leftmargin=2.4em]
  \item \textbf{语音近似}：音节数一致或仅差一音节；
        对应音节的声母、韵母大体对应（细则见下）；
  \item \textbf{场景稳定}：$w'$ 在社群使用中与 $w$
        出现在相同或明确对应的场景
        （如“俺无敌”与“安迪”都出现于奶蛋的“高光”场景）；
  \item \textbf{实例可检索}：存在可公开检索的二创实例
        （视频、帖、弹幕）作为使用证据。
\end{enumerate}

标准 (2)(3) 把“谐音”与“空耳”区分开：
任何人都能给一个词找谐音，但只有社群实际在用的谐音
才是空耳词。\keyw{判定单位是社群行为，不是音系学}。

\subsection{语音距离指标（M 级）}

为了给标准 (1) 一个可排序的度量，本书构造
\keyw{拼音声韵母编辑距离}：

\begin{definition}[声韵母编辑距离（本书构造，\conf{M} 级）]\label{def:08:dist}
  把每个音节表示为有序对（声母 $i$，韵母 $f$），
  忽略声调。两词 $w, w'$ 的距离定义为音节序列上
  最小编辑代价：替换一个声母记 1 分，替换一个韵母记 1 分，
  增删一个音节记 1.5 分；总代价除以较长词的音节数得
  \keyw{归一化距离} $d(w,w') \in [0, +\infty)$。
  $d=0$ 为完全同音（不计声调）。
\end{definition}

三条使用警告。其一，这是\keyw{排序指标}：
它只用于把候选变体按近似程度排序，
\keyw{不}是音系学结论，数值大小没有感知意义。
其二，分值对“声母优先还是韵母优先”等权重设定敏感，
修订者可调整权重，但必须随结果公布所用权重。
其三，$d$ 小不等于等价关系成立——
等价关系由三条标准联合判定，$d$ 只服务标准 (1)。

\begin{example}[演示计算]\label{ex:08:dist}
  以“安迪”（\texttt{an di}）与其三个社群空耳变体为例：
  \begin{enumerate}[label=(\alph*),leftmargin=2.2em]
    \item “俺滴”（\texttt{an di}）：声韵全同，
          $d = 0$；
    \item “安定”（\texttt{an ding}）：音节数相同；
          首音节全同（0 分）；“迪--定”声母同（0 分）、
          韵母 \texttt{i}/\texttt{ing} 替换（1 分）；
          $d = 1/2 = 0.50$；
    \item “俺无敌”（\texttt{an wu di}）：三个音节对两音节，
          取最优对齐“安=俺（0）、迪=敌（0）、增‘无’（1.5）”；
          $d = 1.5/3 = 0.50$。
  \end{enumerate}
  三个变体的 $d$ 均不超过 $0.50$，
  且 (a)(c) 在奶蛋“高光”场景中高频共现、
  (b) 见于“安迪”的普通话化改写场景
  （\conf{C}，第一版教材口径），
  三条标准均满足，故均计入等价项集合。
\end{example}

\begin{nlmethod}[研究方法提示：做一遍的完整步骤]
  修订者复现本节方法的完整步骤：
  (1) 取衍生语清单（表 \ref{tab:06:langs}）；
  (2) 在选定平台以衍生语为关键词检索二创实例，
      记录每条候选空耳的出处链接与访问日期；
  (3) 按 (1)--(3) 三条标准逐项判定，
      记录每条的通过 / 淘汰及理由；
  (4) 对通过项计算 $d$ 并排序；
  (5) 公开判定表（候选、出处、判定、$d$ 值、权重设定）。
  样本规模、时间窗与平台清单必须事先写定——
  这一步骤本身就是第十二章“可复现性清单”的
  一个完整样例：\keyw{方法可复现，结论才可被推翻}。
\end{nlmethod}

\section{AIGC、版权与“变异权”的张力}\label{sec:08:copyright}

AIGC 把第三章的版权张力推向了一个新形态。
正典时代的侵权是\keyw{复制}（盗图、盗播），
AIGC 时代的侵权是\keyw{变异}（AI 重绘、批量生成）——
后者在法律上更难定性：变异产物与原形象的相似度
落在哪个区间构成侵权，目前没有稳定判例口径，
本书不给出法律意见（边界声明见表 \ref{tab:02:boundary}）。
可以指认的只有行为层面的结构：

\begin{enumerate}[label=\arabic*.]
  \item \textbf{生成即上游}：奶蛙案例显示，变异形象可以在
        平台侧“无主”地诞生（智能体头像），
        再由社群认领——侵权行为与侵权意图在起点处脱钩；
  \item \textbf{批量即规模}：单个变体的法律意义有限，
        权利人实际面对的是批量生成造成的
        \keyw{执法对象海量化}；
  \item \textbf{变异即传播}：官方对“丑化二创”的压制
        （第四章 \ref{sec:04:official}）恰好撞上
        M5 机制（\ref{sec:07:reverse}）——
        AIGC 供给的无限性与压制的有限性相遇，
        构成第九章“迷因生态缝隙”的生产端基础。
\end{enumerate}

\begin{nlcase}[案例 8-1：从“变形”到“变异”的权利视角]
  对比两个场景。场景一：某商家未经授权印刷奶龙图案销售
  ——经典侵权，权利边界清晰，维权路径成熟
  （快速维权机制，\conf{A}\cite{jiang2026nailong}）。
  场景二：某用户用图像模型把奶龙“变”成奶蛙并发布
  ——变异产物与原形象的相似度可辩、
  生成过程无“复制”动作、发布者与传播链解耦。
  两个场景在权利人眼中同属“形象失真”的损害
  （谢承恩列举的第一种侵权形态即“IP 形象严重失真”，
  \conf{A}\cite{jiang2026nailong}），
  但在现行工具与平台结构下，
  第二场景的可诉性、可执行性都显著更低。
  这一落差解释了为什么官方反应集中在
  \keyw{平台行为}（关评、拉黑账号）而非
  \keyw{图像本体}（要求删除变异形象）——
  平台行为是可执行的，图像本体不是。
\end{nlcase}

\section{本章小结}\label{sec:08:summary}

本章拆解了奶蛙二创的五环节生产链，
指出变异引擎的两个时点（2024 星野智能体、
2026 视频生成模型，均 \conf{C} 级），
并给出形式稳定性的经济学版本
（命题 \ref{prop:08:asym}：内容变异成本趋零、
形式构件模板固化，\conf{M} 级）。
方法论上，本章兑现了全书的可复现承诺：
空耳词等价关系的三条判定标准与
声韵母编辑距离（定义 \ref{def:08:dist}，
M 级排序指标），并给出完整复现步骤。
张力层面，AIGC 把侵权从“复制”变为“变异”，
可执行的对象从图像本体移向平台行为——
这是第九章讨论官方与二创关系的生产端前提。
遗留问题：生产端解释了“供给为何无限”，
但没有回答“需求为何存在”——
那需要回到第七章的机制假设与第十一章的受体研究；
同时，供给泛滥的外部性如何被治理
（平台、权利人、社群三方），
是第九章产业分析与第十章校园协商的共同背景。

\begin{nlexercises}
  \item 用定义 \ref{def:08:dist} 计算
        “嘎嘎滴辣虾”（\texttt{ga ga di la xia}）与
        “刚刚点了谁”（\texttt{gang gang dian le shui}）
        的归一化距离，写出你的对齐方案。
        再解释：为什么即使 $d$ 较大，
        该空耳仍可成立（提示：标准 (2)(3)）。
  \item 命题 \ref{prop:08:asym} 的推论二预测
        “改结尾套路”的二创传播表现更差。
        设计一个最小对照：
        选什么样的样本、编码哪两个变量、
        用什么统计量比较，并说明
        如何排除“改结尾的作者本来就不熟练”这一混杂因素。
  \item “生产成本由个体承担、拥挤成本由公共信息流承担”
        是一个外部性结构。仿照排放税的思路，
        提出一个平台侧的治理机制设计
        （如对高频发布者的某个约束），
        并分析它会触发第七章哪个机制的反弹
        （提示：M5）。
\end{nlexercises}

\nlkeywords AIGC、生产链、变异引擎、Seedance、成本不对称、
模板固化、空耳（soramimi）、声韵母编辑距离、
判定标准、低质内容、外部性、变异权
